% =====================================================================
% 第 3 章 测试与结果（owner: writer-results）
%
% 数字来源约定（均为本轮 2026-10-20 重新构建/重新运行的结果）：
%   % src: tasks F##  —— reports/data/results-tasks.md（任务 1–3 实测，F1–F42）
%   % src: verify F## —— reports/data/results-verify.md（测试/交叉验证/Bonus，F1–F18）
%   % src: pre P##    —— reports/data/results-preprocess.md（预处理复核，P1–P6）
% 本章不使用 README.md/notes.md/plan.md 里的历史数字；开发期 bug 只作定性叙述，
% 修复后的验证结论一律取本轮实测结果。每个表格/段落前标注其证据编号。
% =====================================================================

\newpage
\section{测试与结果分析}
\label{sec:results}

% 本章有 14 张表、3 张图，局部收紧缩放间距，减少浮动体造成的整页空白；
% 章末恢复 article 默认值，避免影响后续章节。
\setlength{\textfloatsep}{8pt plus 2pt minus 2pt}
\setlength{\floatsep}{8pt plus 2pt minus 2pt}
\setlength{\intextsep}{8pt plus 2pt minus 2pt}

\subsection{词频统计与停用词选择}

\subsubsection{词频统计结果}
\label{sec:wordcount}

默认配置只统计正文行。195 个文档共切分出 866\,626 个词元，其中不同词项
25\,705 个，仅出现一次的词项 9\,753 个。核心统计量见表~\ref{tab:wordcount}。

\begin{table}[!htbp]
  \centering
  \caption{默认配置下的词频统计结果}
  \label{tab:wordcount}
  \small
  \begin{tabular}{lr}
    \toprule
    指标 & 数值 \\
    \midrule
    文档数 & 195 \\
    正文词元数 & 866\,626 \\
    不同词项数 & 25\,705 \\
    仅出现一次的词项数 & 9\,753 \\
    \bottomrule
  \end{tabular}
\end{table}

默认配置下词频最高的 10 个词项见表~\ref{tab:top10}。
% src: tasks F7
\begin{table}[!htbp]
  \centering
  \caption{默认配置下词频最高的 10 个词项（含文档频率与两种占比）}
  \label{tab:top10}
  \footnotesize
  \renewcommand{\arraystretch}{1.0}
  \begin{tabular}{r l r r r r}
    \toprule
    排名 & 词项 & 词频 & 文档频率 & 文档频率比值 & 词频占比 \\
    \midrule
    1 & \code{the} & 28\,436 & 177 & 0.9077 & 0.032812 \\
    2 & \code{and} & 25\,553 & 189 & 0.9692 & 0.029486 \\
    3 & \code{i}   & 21\,131 & 154 & 0.7897 & 0.024383 \\
    4 & \code{to}  & 20\,302 & 187 & 0.9590 & 0.023426 \\
    5 & \code{of}  & 16\,976 & 175 & 0.8974 & 0.019589 \\
    6 & \code{a}   & 14\,527 & 138 & 0.7077 & 0.016763 \\
    7 & \code{you} & 13\,998 & 75  & 0.3846 & 0.016152 \\
    8 & \code{my}  & 12\,934 & 159 & 0.8154 & 0.014925 \\
    9 & \code{that}& 11\,513 & 172 & 0.8821 & 0.013285 \\
    10& \code{in}  & 11\,479 & 181 & 0.9282 & 0.013246 \\
    \bottomrule
  \end{tabular}
  \par\smallskip
  \footnotesize 注：文档频率比值为文档频率/195，词频占比为词频/866\,626；数值取自  % src: tasks F7
  \code{wordcount-default.csv} 的权威副本。
\end{table}

% figures/zipf.tex -- 词频—排名分布（双对数），Zipf 定律
%
% 数据：reports/data/csv/tasks/wordcount-default.csv（本轮 mse count，默认管线）
% 由 figures/data/make_figures_data.py 生成 figures/data/zipf.dat 与 zipf-meta.tex。
% 主文档只写 % figures/zipf.tex -- 词频—排名分布（双对数），Zipf 定律
%
% 数据：reports/data/csv/tasks/wordcount-default.csv（本轮 mse count，默认管线）
% 由 figures/data/make_figures_data.py 生成 figures/data/zipf.dat 与 zipf-meta.tex。
% 主文档只写 % figures/zipf.tex -- 词频—排名分布（双对数），Zipf 定律
%
% 数据：reports/data/csv/tasks/wordcount-default.csv（本轮 mse count，默认管线）
% 由 figures/data/make_figures_data.py 生成 figures/data/zipf.dat 与 zipf-meta.tex。
% 主文档只写 \input{figures/zipf}。
\input{figures/data/zipf-meta}
\begin{figure}[!htbp]
\centering
\begin{tikzpicture}
\begin{axis}[
  width=0.94\linewidth, height=6.0cm,
  xmode=log, ymode=log,
  xmin=1, xmax=40000, ymin=1, ymax=1e5,
  xlabel={rank (log scale)}, ylabel={term frequency tf (log scale)},
  xtick={1,10,100,1000,10000},
  ytick={1,10,100,1000,10000},
  tick label style={font=\scriptsize},
  label style={font=\small},
  grid=major, grid style={black!8},
  legend style={draw=none, fill=none, font=\scriptsize,
                at={(0.98,0.98)}, anchor=north east},
  legend cell align=left,
]
  % 实测 tf(rank)
  \addplot[only marks, mark=*, mark size=1.0pt, color=mseblue]
    table[x=rank, y=tf] {figures/data/zipf.dat};
  \addlegendentry{measured tf (\ZipfTerms{} terms)}
  % 1/rank 参考线（按 rank 1 处缩放对齐）
  \addplot[dashed, line width=0.8pt, color=msegray]
    table[x=rank, y=ref] {figures/data/zipf.dat};
  \addlegendentry{$1/\mathrm{rank}$ reference (scaled at rank 1)}
  % 最频词标注
  \node[font=\scriptsize, text=msegray, anchor=west] (topword)
    at (axis cs:1.15,42000) {the};
  \draw[->, draw=msegray, line width=0.4pt] (topword.west) -- (axis cs:1.0,3.1e4);
  % hapax 尾标注
  \node[font=\scriptsize, text=msegray, align=left, anchor=south west]
    at (rel axis cs:0.02,0.03)
    {hapax tail: \ZipfHapax{} terms (\ZipfHapaxPct\% of vocabulary)\\
     occur exactly once (tf $=1$)};
\end{axis}
\end{tikzpicture}
\caption{莎士比亚语料的词频—排名分布（双对数）：实测点近似一条直线，符合
Zipf 定律；\ZipfTerms{} 个词中有 \ZipfHapax{} 个只出现一次（占词表
\ZipfHapaxPct\%），构成需要倒排索引兜住的长尾。}%
\label{fig:zipf}
\end{figure}
。
% zipf-meta.tex -- generated by figures/data/make_figures_data.py
\newcommand{\ZipfTerms}{25\,705}
\newcommand{\ZipfHapax}{9\,753}
\newcommand{\ZipfHapaxPct}{37.9}

\begin{figure}[!htbp]
\centering
\begin{tikzpicture}
\begin{axis}[
  width=0.94\linewidth, height=6.0cm,
  xmode=log, ymode=log,
  xmin=1, xmax=40000, ymin=1, ymax=1e5,
  xlabel={rank (log scale)}, ylabel={term frequency tf (log scale)},
  xtick={1,10,100,1000,10000},
  ytick={1,10,100,1000,10000},
  tick label style={font=\scriptsize},
  label style={font=\small},
  grid=major, grid style={black!8},
  legend style={draw=none, fill=none, font=\scriptsize,
                at={(0.98,0.98)}, anchor=north east},
  legend cell align=left,
]
  % 实测 tf(rank)
  \addplot[only marks, mark=*, mark size=1.0pt, color=mseblue]
    table[x=rank, y=tf] {figures/data/zipf.dat};
  \addlegendentry{measured tf (\ZipfTerms{} terms)}
  % 1/rank 参考线（按 rank 1 处缩放对齐）
  \addplot[dashed, line width=0.8pt, color=msegray]
    table[x=rank, y=ref] {figures/data/zipf.dat};
  \addlegendentry{$1/\mathrm{rank}$ reference (scaled at rank 1)}
  % 最频词标注
  \node[font=\scriptsize, text=msegray, anchor=west] (topword)
    at (axis cs:1.15,42000) {the};
  \draw[->, draw=msegray, line width=0.4pt] (topword.west) -- (axis cs:1.0,3.1e4);
  % hapax 尾标注
  \node[font=\scriptsize, text=msegray, align=left, anchor=south west]
    at (rel axis cs:0.02,0.03)
    {hapax tail: \ZipfHapax{} terms (\ZipfHapaxPct\% of vocabulary)\\
     occur exactly once (tf $=1$)};
\end{axis}
\end{tikzpicture}
\caption{莎士比亚语料的词频—排名分布（双对数）：实测点近似一条直线，符合
Zipf 定律；\ZipfTerms{} 个词中有 \ZipfHapax{} 个只出现一次（占词表
\ZipfHapaxPct\%），构成需要倒排索引兜住的长尾。}%
\label{fig:zipf}
\end{figure}
。
% zipf-meta.tex -- generated by figures/data/make_figures_data.py
\newcommand{\ZipfTerms}{25\,705}
\newcommand{\ZipfHapax}{9\,753}
\newcommand{\ZipfHapaxPct}{37.9}

\begin{figure}[!htbp]
\centering
\begin{tikzpicture}
\begin{axis}[
  width=0.94\linewidth, height=6.0cm,
  xmode=log, ymode=log,
  xmin=1, xmax=40000, ymin=1, ymax=1e5,
  xlabel={rank (log scale)}, ylabel={term frequency tf (log scale)},
  xtick={1,10,100,1000,10000},
  ytick={1,10,100,1000,10000},
  tick label style={font=\scriptsize},
  label style={font=\small},
  grid=major, grid style={black!8},
  legend style={draw=none, fill=none, font=\scriptsize,
                at={(0.98,0.98)}, anchor=north east},
  legend cell align=left,
]
  % 实测 tf(rank)
  \addplot[only marks, mark=*, mark size=1.0pt, color=mseblue]
    table[x=rank, y=tf] {figures/data/zipf.dat};
  \addlegendentry{measured tf (\ZipfTerms{} terms)}
  % 1/rank 参考线（按 rank 1 处缩放对齐）
  \addplot[dashed, line width=0.8pt, color=msegray]
    table[x=rank, y=ref] {figures/data/zipf.dat};
  \addlegendentry{$1/\mathrm{rank}$ reference (scaled at rank 1)}
  % 最频词标注
  \node[font=\scriptsize, text=msegray, anchor=west] (topword)
    at (axis cs:1.15,42000) {the};
  \draw[->, draw=msegray, line width=0.4pt] (topword.west) -- (axis cs:1.0,3.1e4);
  % hapax 尾标注
  \node[font=\scriptsize, text=msegray, align=left, anchor=south west]
    at (rel axis cs:0.02,0.03)
    {hapax tail: \ZipfHapax{} terms (\ZipfHapaxPct\% of vocabulary)\\
     occur exactly once (tf $=1$)};
\end{axis}
\end{tikzpicture}
\caption{莎士比亚语料的词频—排名分布（双对数）：实测点近似一条直线，符合
Zipf 定律；\ZipfTerms{} 个词中有 \ZipfHapax{} 个只出现一次（占词表
\ZipfHapaxPct\%），构成需要倒排索引兜住的长尾。}%
\label{fig:zipf}
\end{figure}


前十个词几乎全是功能词，合计占 20.41\% 的词元。图~\ref{fig:zipf} 的双对数曲线
进一步显示，高频词集中在曲线左端，绝大多数词项位于右端的低频长尾。仅出现一次的词
占到全部词项的 37.94\%，却只覆盖 1.125\% 的词元，说明语料具有明显的“短词表头部、
长低频尾部”特征。由此不能仅凭词频高低判断停用词，下一步还需要比较文档覆盖范围和
文档间分布。

\subsubsection{停用词选择结果}
\label{sec:stopword-results}

% src: tasks F10、F11
表~\ref{tab:candidates} 是候选规则对文档频率阈值的敏感性扫描：固定词频占比
$\ge 0.001$，把文档频率比值阈值从 0.90 降到 0.40。
\begin{table}[!htbp]
  \centering
  \caption{停用词候选带对文档频率比值阈值的敏感性（固定词频占比 $\ge 0.001$）}
  \label{tab:candidates}
  \footnotesize
  \renewcommand{\arraystretch}{1.0}
  \begin{tabularx}{\linewidth}{c c c X}
    \toprule
    文档频率比值阈值 $T$ & 候选数 & 相对上一档新增 & 新增词 \\
    \midrule
    0.90 & 4  & --- & \code{the and to in} \\
    0.75 & 9  & +5  & \code{i of my that but} \\
    0.70 & 14 & +5  & \code{a is not for with} \\
    0.65 & 16 & +2  & \code{so love} \\
    0.60 & 21 & +5  & \code{me be thou thy thee} \\
    0.50 & 31 & +10 & \code{it this as by all do on from which when} \\
    0.40 & 48 & +17 & \code{his have what no if are o or more then their hath than yet mine time doth} \\
    \bottomrule
  \end{tabularx}
  \par\smallskip
  \footnotesize 注：规则为文档频率比值 $\ge T$ 且词频占比 $\ge 0.001$；静态 39 词表就是  % src: tasks F11
  在 0.40 档 48 词候选带上人工取舍的结果。
\end{table}

0.75 档只有 9 个纯功能词；0.65 档内容词 \code{love} 进入候选（文档频率比值 0.651）；0.40
档就混入 \code{time}、\code{thou/thee} 等不该删除的词。纯统计规则能可靠框出"最像
停用词"的一小段，却无法自己决定线的位置——静态表是在 48 词候选带上人工保留结构性
功能词的结果，下一节用速率熵自动完成同一步。

\subsubsection{自动推导结果}

两特征候选带把最后一步留给人工审核；加上"平"（速率熵）后，规则能否自动得到与人工表
几乎相同的停用词？表~\ref{tab:stopword-sweep} 是四组阈值与 39 词静态表的双向差集。
% src: tasks F13、F14
\begin{table}[!htbp]
  \centering
  \caption{三特征停用词推导的阈值扫描：与静态 39 词表的差异}
  \label{tab:stopword-sweep}
  \footnotesize
  \renewcommand{\arraystretch}{1.0}
  \begin{tabularx}{\linewidth}{c c c c X}
    \toprule
    (文档频率比值, 速率熵) & 推导词数 & 重叠 & 静态表独有 & 推导表独有 \\
    \midrule
    (0.75, 0.93) & 9  & 9  & 30 & --- \\
    (0.50, 0.95) & 23 & 23 & 16 & --- \\
    (0.40, 0.93) & 39 & 35 & \code{more or than their} & \code{hath o thee thou} \\
    (0.30, 0.93) & 50 & 35 & \code{more or than their} & \code{at hath he him how let now o thee there thou was well would you} \\
    \bottomrule
  \end{tabularx}
  \par\smallskip
  \footnotesize 注：词频占比固定为 $\ge 0.001$，文档频率比值取第一列；  % src: tasks F13
  (0.75,0.93) 档 9 词与两特征候选带一致，(0.40,0.93) 档复现静态表所在的 39 词带。
\end{table}

(0.40, 0.93) 档推导出 39 词、与静态表重叠 35 个；差异恰好是 4 对可解释的词：自动
推导多收 \code{hath o thee thou} 四个古体功能词，静态表多收 \code{more or than
their} 四个兼作比较、并列用的词。前者在原文本中同样"宽而平"，后者是人工按语义判断的
取舍。也就是说，三特征规则把"是不是停用词"从人工经验变成可复算的判定，结果与人工表
只差 4 个词且差异可逐词解释；换一档阈值只改变严格度（9/23/39/50 词），不改变规则
结构。四个输出文件都带出处头且不含时间戳，换参数可以直接重跑比较。

图~\ref{fig:entropy} 把功能词、内容词与专名的速率熵画成条形图并标出 0.93 阈值线。

% figures/entropy.tex -- 速率熵（entropy_rate_norm）判别停用词
%
% 数据：reports/data/csv/tasks/wordcount-default.csv（本轮 mse count，默认管线）
% 由 figures/data/make_figures_data.py 生成 figures/data/entropy-*.dat。
% 主文档只写 % figures/entropy.tex -- 速率熵（entropy_rate_norm）判别停用词
%
% 数据：reports/data/csv/tasks/wordcount-default.csv（本轮 mse count，默认管线）
% 由 figures/data/make_figures_data.py 生成 figures/data/entropy-*.dat。
% 主文档只写 % figures/entropy.tex -- 速率熵（entropy_rate_norm）判别停用词
%
% 数据：reports/data/csv/tasks/wordcount-default.csv（本轮 mse count，默认管线）
% 由 figures/data/make_figures_data.py 生成 figures/data/entropy-*.dat。
% 主文档只写 \input{figures/entropy}。
\begin{figure}[!htbp]
\centering
\begin{tikzpicture}
\begin{axis}[
  width=0.96\linewidth, height=5.3cm,
  ymin=0, ymax=1.06,
  ytick={0,0.2,0.4,0.6,0.8,1.0},
  ylabel={entropy\_rate\_norm},
  symbolic x coords={the,and,i,to,of,is,it,that,love,king,heart,sweet,man,osric,banquo},
  xtick={the,and,i,to,of,is,it,that,love,king,heart,sweet,man,osric,banquo},
  xticklabel style={font=\scriptsize, rotate=40, anchor=north east, inner sep=1pt},
  tick label style={font=\scriptsize},
  label style={font=\small},
  enlarge x limits=0.04,
  grid=major, grid style={black!8},
  legend columns=4,
  legend style={at={(0.5,-0.23)}, anchor=north, draw=none, font=\scriptsize},
  legend cell align=left,
]
  % 功能词：速率熵接近 1
  \addplot[ycomb, color=mseblue, line width=1.0pt,
           mark=*, mark size=1.8pt, mark options={fill=mseblue}]
    table[x=term, y=rate] {figures/data/entropy-function.dat};
  \addlegendentry{function words}
  % 内容词：明显偏低
  \addplot[ycomb, color=msegray, line width=0.9pt,
           mark=square*, mark size=1.6pt, mark options={fill=msegray}]
    table[x=term, y=rate] {figures/data/entropy-content.dat};
  \addlegendentry{content words}
  % 专名：df=1 时按定义为 0
  \addplot[ycomb, color=black!35, line width=0.9pt,
           mark=triangle*, mark size=2.0pt, mark options={fill=black!35}]
    table[x=term, y=rate] {figures/data/entropy-name.dat};
  \addlegendentry{proper nouns (df $=1$)}
  % 规则阈值 0.93
  \addplot[dashed, line width=0.8pt, color=black!65, mark=none]
    table[x=term, y=threshold] {figures/data/entropy-threshold.dat};
  \addlegendentry{0.93 threshold}
\end{axis}
\end{tikzpicture}
\caption{速率熵判别：功能词全部贴近 1（均在 0.96 以上），专名（df${}=1$）
按定义为 0，内容词落在 0.83–0.94 的低区间；0.93 一条横线即可把功能词与其余
词分开，个别越线的内容词再被 df 比例门（$\ge 0.4$）挡住，说明该特征能替代
人工逐个审核。}%
\label{fig:entropy}
\end{figure}
。
\begin{figure}[!htbp]
\centering
\begin{tikzpicture}
\begin{axis}[
  width=0.96\linewidth, height=5.3cm,
  ymin=0, ymax=1.06,
  ytick={0,0.2,0.4,0.6,0.8,1.0},
  ylabel={entropy\_rate\_norm},
  symbolic x coords={the,and,i,to,of,is,it,that,love,king,heart,sweet,man,osric,banquo},
  xtick={the,and,i,to,of,is,it,that,love,king,heart,sweet,man,osric,banquo},
  xticklabel style={font=\scriptsize, rotate=40, anchor=north east, inner sep=1pt},
  tick label style={font=\scriptsize},
  label style={font=\small},
  enlarge x limits=0.04,
  grid=major, grid style={black!8},
  legend columns=4,
  legend style={at={(0.5,-0.23)}, anchor=north, draw=none, font=\scriptsize},
  legend cell align=left,
]
  % 功能词：速率熵接近 1
  \addplot[ycomb, color=mseblue, line width=1.0pt,
           mark=*, mark size=1.8pt, mark options={fill=mseblue}]
    table[x=term, y=rate] {figures/data/entropy-function.dat};
  \addlegendentry{function words}
  % 内容词：明显偏低
  \addplot[ycomb, color=msegray, line width=0.9pt,
           mark=square*, mark size=1.6pt, mark options={fill=msegray}]
    table[x=term, y=rate] {figures/data/entropy-content.dat};
  \addlegendentry{content words}
  % 专名：df=1 时按定义为 0
  \addplot[ycomb, color=black!35, line width=0.9pt,
           mark=triangle*, mark size=2.0pt, mark options={fill=black!35}]
    table[x=term, y=rate] {figures/data/entropy-name.dat};
  \addlegendentry{proper nouns (df $=1$)}
  % 规则阈值 0.93
  \addplot[dashed, line width=0.8pt, color=black!65, mark=none]
    table[x=term, y=threshold] {figures/data/entropy-threshold.dat};
  \addlegendentry{0.93 threshold}
\end{axis}
\end{tikzpicture}
\caption{速率熵判别：功能词全部贴近 1（均在 0.96 以上），专名（df${}=1$）
按定义为 0，内容词落在 0.83–0.94 的低区间；0.93 一条横线即可把功能词与其余
词分开，个别越线的内容词再被 df 比例门（$\ge 0.4$）挡住，说明该特征能替代
人工逐个审核。}%
\label{fig:entropy}
\end{figure}
。
\begin{figure}[!htbp]
\centering
\begin{tikzpicture}
\begin{axis}[
  width=0.96\linewidth, height=5.3cm,
  ymin=0, ymax=1.06,
  ytick={0,0.2,0.4,0.6,0.8,1.0},
  ylabel={entropy\_rate\_norm},
  symbolic x coords={the,and,i,to,of,is,it,that,love,king,heart,sweet,man,osric,banquo},
  xtick={the,and,i,to,of,is,it,that,love,king,heart,sweet,man,osric,banquo},
  xticklabel style={font=\scriptsize, rotate=40, anchor=north east, inner sep=1pt},
  tick label style={font=\scriptsize},
  label style={font=\small},
  enlarge x limits=0.04,
  grid=major, grid style={black!8},
  legend columns=4,
  legend style={at={(0.5,-0.23)}, anchor=north, draw=none, font=\scriptsize},
  legend cell align=left,
]
  % 功能词：速率熵接近 1
  \addplot[ycomb, color=mseblue, line width=1.0pt,
           mark=*, mark size=1.8pt, mark options={fill=mseblue}]
    table[x=term, y=rate] {figures/data/entropy-function.dat};
  \addlegendentry{function words}
  % 内容词：明显偏低
  \addplot[ycomb, color=msegray, line width=0.9pt,
           mark=square*, mark size=1.6pt, mark options={fill=msegray}]
    table[x=term, y=rate] {figures/data/entropy-content.dat};
  \addlegendentry{content words}
  % 专名：df=1 时按定义为 0
  \addplot[ycomb, color=black!35, line width=0.9pt,
           mark=triangle*, mark size=2.0pt, mark options={fill=black!35}]
    table[x=term, y=rate] {figures/data/entropy-name.dat};
  \addlegendentry{proper nouns (df $=1$)}
  % 规则阈值 0.93
  \addplot[dashed, line width=0.8pt, color=black!65, mark=none]
    table[x=term, y=threshold] {figures/data/entropy-threshold.dat};
  \addlegendentry{0.93 threshold}
\end{axis}
\end{tikzpicture}
\caption{速率熵判别：功能词全部贴近 1（均在 0.96 以上），专名（df${}=1$）
按定义为 0，内容词落在 0.83–0.94 的低区间；0.93 一条横线即可把功能词与其余
词分开，个别越线的内容词再被 df 比例门（$\ge 0.4$）挡住，说明该特征能替代
人工逐个审核。}%
\label{fig:entropy}
\end{figure}


功能词 the、and、i、to、of、is、it、that 的速率熵都在 0.95 以上（the 0.9748、
and 0.9789），内容词 love、king、heart、sweet 落在 0.82--0.93（0.9271、0.8260、
0.8313、0.8925），只出现在一篇文档的专名 osric、banquo 按定义为 0。0.93 的线正好切在
功能词与内容词之间；\code{man} 的速率熵 0.9388 略高于阈值线，但文档频率比值只有 0.2359，
被"宽"这一条件挡在候选带外——三个条件是"与"关系，速率熵补上的正是前两个特征缺失的
"分布是否均匀"这一维。

最终停用词的选择是 \code{the and i to of a my that in is not me it for with be his this but have as so what by no all do if are on from or which more then their when than yet}。

\subsection{倒排索引与查询结果}

\subsubsection{倒排索引构建结果}
\label{sec:index-build}

正式倒排索引采用 39 词停用表并执行词干化。构建完成后的结果见
表~\ref{tab:index-configs}。
% src: tasks F18、F21
\begin{table}[!htbp]
  \centering
  \caption{正式倒排索引的最终构建结果}
  \label{tab:index-configs}
  \small
  \begin{tabular}{lr}
    \toprule
    指标 & 结果 \\
    \midrule
    文档数 & 195 \\
    保留词元数 & 556\,866 \\
    词项数 & 14\,937 \\
    倒排记录数 & 117\,704 \\
    位置数 & 556\,866 \\
    第一版索引大小 & 3\,685\,655 字节 \\
    第二版索引大小 & 1\,725\,694 字节 \\
    构建时间 & 0.19\,秒 \\
    \bottomrule
  \end{tabular}
\end{table}

表~\ref{tab:dict-compare} 比较三种词典实现构建同一份正式索引时的代价。
% src: tasks F20
\begin{table}[!htbp]
  \centering
  \caption{三种词典实现的构建代价（正式配置，各构建 3 次）}
  \label{tab:dict-compare}
  \footnotesize
  \renewcommand{\arraystretch}{1.0}
  \begin{tabular}{l r r r l}
    \toprule
    词典 & 比较次数 & 旋转次数 & 再哈希次数 & 构建时间（内建 / 外部中位） \\
    \midrule
    伸展树（主） & 6\,392\,724 & 5\,850\,795 & 0 & 0.19\,s / 0.2610\,s \\
    哈希表 & 688\,494 & 0 & 11 & 0.09\,s / \textbf{0.1603\,s} \\
    \code{std::map}（对照） & 9\,435\,694 & 0 & 0 & 0.14\,s / 0.2084\,s \\
    \bottomrule
  \end{tabular}
  \par\smallskip
  \footnotesize 注：比较次数、旋转次数、再哈希次数为 C++ 内建计数；外部墙钟含进程  % src: tasks F20
  启动与文件读写，受并行负载影响只作量级比较。
\end{table}

从构建速度看，哈希表最快，标准映射次之，伸展树最慢；从字符串比较次数看，哈希表也
明显少于两种树结构。伸展树的主要代价是 585 万次旋转，但它的中序遍历可以直接得到
有序词项表。项目因此以伸展树作为主词典，以哈希表作为速度基线，以标准映射作为平衡树
对照。第二版索引在第一版基础上降至 172 万字节，约为原来的 46.82\%。

\subsubsection{查询结果}
\label{sec:query}

本节使用正式索引测试单词、短语和布尔查询，并检查无结果时能否说明原因。
代表性查询见表~\ref{tab:query-examples}。
% src: tasks F27–F33
\begin{table}[!htbp]
  \centering
  \caption{正式索引上的代表性查询结果}
  \label{tab:query-examples}
  \footnotesize
  \renewcommand{\arraystretch}{1.0}
  \begin{tabularx}{\linewidth}{l l r X}
    \toprule
    查询 & 类型 & 文档数 & 结果说明 \\
    \midrule
    \code{osric} & 单词 & 1 & 文档 30《Hamlet》，可定位到原句 \\
    \code{love} & 单词 & 138 & 高频内容词，返回全部包含该词的文档 \\
    \code{good night} & 短语 & 23 & 两个词必须位置相邻 \\
    \code{king richard} & 短语 & 7 & 两个词均存在并通过位置验证 \\
    \code{my lord} & 短语 & 42 & \code{my} 被停用词过滤，实际只查 \code{lord} \\
    \code{to be or not to be} & 短语 & 0 & 全部由停用词组成 \\
    \code{king queen} & 布尔 & 33 & 默认要求两个词全部命中 \\
    \code{the} & 单词 & 0 & 停用词未进入索引 \\
    \code{notaword} & 单词 & 0 & 索引词典中没有该词 \\
    \bottomrule
  \end{tabularx}
\end{table}

单词查询直接返回包含目标词项的文档编号和标题。结果按命中词数、词频总和和文档编号
稳定排序。加上 \code{--pretty} 后，程序会重新运行文本处理管线，在结果中附上词语
所在的原句，便于核对命中位置。

短语查询要验证多个词项在过滤停用词后的词元序列中连续出现，不能只求两个倒排记录表
的交集。\code{good night} 命中 23 篇，只有两个词完整相邻出现的文档会被保留。由于
位置定义在过滤后的词元序列上，\code{my lord} 中的 \code{my} 不占位置，系统实际执行
的是 \code{lord} 的单词查询。全部由停用词组成的短语没有可定位的词项，因此返回空结果。

布尔查询默认要求所有查询词都命中。\code{king queen} 因此返回 33 篇；使用
\code{--min-match 1} 后只需命中任意一个词，结果增加到 43 篇。无结果时，程序会区分
“停用词未进入索引”和“索引词典中没有该词”，避免把所有空结果混为一类。

\subsection{阈值实验与复杂度分析}

\subsubsection{阈值实验结果}
\label{sec:threshold}

四个参数都只作用于查询。全部关闭时，同一查询的结果作为基线。实验每次只改变一个参数，
共选 15 条查询，覆盖高频词、中频词、低频词、专名、普通短语、含停用词短语、布尔查询
和停用词。对全部查询测试 5 档最大文档频率、4 档最小文档频率和 3 档结果截断；最少
命中词数只适用于 king queen 与 sweet sword moon 两条布尔查询，取 3 档。由此得到
$15+15\times(5+4+3)+2\times3=201$ 条测量。每条记录返回数、相对基线的 Jaccard
相似度、说明信息条数和耗时，全部使用同一份正式索引。图~\ref{fig:threshold} 展示了
四个参数对返回文档数的影响。

% figures/threshold.tex -- 四个查询阈值对返回文档数的影响（2x2 面板）
%
% 数据：reports/data/csv/tasks/threshold_matrix.csv（本轮 tools/threshold_matrix.py
% 的 201 条测量；off 行即无阈值基线）。由 figures/data/make_figures_data.py
% 生成 figures/data/threshold-*.dat。主文档只写 % figures/threshold.tex -- 四个查询阈值对返回文档数的影响（2x2 面板）
%
% 数据：reports/data/csv/tasks/threshold_matrix.csv（本轮 tools/threshold_matrix.py
% 的 201 条测量；off 行即无阈值基线）。由 figures/data/make_figures_data.py
% 生成 figures/data/threshold-*.dat。主文档只写 % figures/threshold.tex -- 四个查询阈值对返回文档数的影响（2x2 面板）
%
% 数据：reports/data/csv/tasks/threshold_matrix.csv（本轮 tools/threshold_matrix.py
% 的 201 条测量；off 行即无阈值基线）。由 figures/data/make_figures_data.py
% 生成 figures/data/threshold-*.dat。主文档只写 \input{figures/threshold}。
\pgfplotsset{
  % 面板公共设置
  figthaxis/.style={
    width=7.10cm, height=2.90cm,
    ymin=0,
    xtick style={draw=black!40},
    tick label style={font=\scriptsize},
    label style={font=\small},
    title style={font=\small},
    grid=major, grid style={black!8},
    unbounded coords=discard,
  },
  % 每条查询一条曲线：颜色/线型/标记三者组合区分
  figthlove/.style={color=mseblue, line width=0.9pt, mark=*, mark size=1.1pt},
  figthsword/.style={color=mseblue, dashed, line width=0.9pt,
                     mark=square*, mark size=1.0pt},
  figthphrase/.style={color=msegray, line width=0.9pt,
                      mark=triangle*, mark size=1.3pt},
  figthkq/.style={color=msegray, dashed, line width=0.9pt,
                  mark=diamond*, mark size=1.1pt},
  figthssm/.style={color=black!65, densely dotted, line width=1.0pt,
                   mark=o, mark size=1.1pt},
}
\begin{figure}[!htbp]
\centering
\captionsetup[subfigure]{font=scriptsize,skip=2pt}
% ---------------- (a) max-df ----------------
\begin{subfigure}{0.49\linewidth}
  \centering
  \begin{tikzpicture}
  \begin{axis}[figthaxis,
    xlabel={max-df threshold}, ylabel={result documents},
    xtick={0,1,2,3,4,5}, xticklabels={0.1,0.3,0.5,0.7,0.9,off},
    xmin=-0.55, xmax=5.55,
    legend to name={figthresholdlegend}, legend columns=3,
    legend style={draw=none, font=\scriptsize},
  ]
    \addplot[figthlove] table[x=idx, y=love] {figures/data/threshold-max-df.dat};
    \addlegendentry{love (word)}
    \addplot[figthsword] table[x=idx, y=sword] {figures/data/threshold-max-df.dat};
    \addlegendentry{sword (word)}
    \addplot[figthphrase] table[x=idx, y=kr_phrase] {figures/data/threshold-max-df.dat};
    \addlegendentry{king richard (phrase)}
    \addplot[figthkq] table[x=idx, y=kq_bool] {figures/data/threshold-max-df.dat};
    \addlegendentry{king queen (bool)}
    \addplot[figthssm] table[x=idx, y=ssm_bool] {figures/data/threshold-max-df.dat};
    \addlegendentry{sweet sword moon (bool)}
  \end{axis}
  \end{tikzpicture}
  \caption{max-df：裁掉 df 比例过高的查询词}
\end{subfigure}\hfill
% ---------------- (b) min-df ----------------
\begin{subfigure}{0.49\linewidth}
  \centering
  \begin{tikzpicture}
  \begin{axis}[figthaxis,
    xlabel={min-df threshold},
    xtick={0,1,2,3,4}, xticklabels={off,2,5,20,50},
    xmin=-0.55, xmax=4.55,
  ]
    \addplot[figthlove, forget plot] table[x=idx, y=love] {figures/data/threshold-min-df.dat};
    \addplot[figthsword, forget plot] table[x=idx, y=sword] {figures/data/threshold-min-df.dat};
    \addplot[figthphrase, forget plot] table[x=idx, y=kr_phrase] {figures/data/threshold-min-df.dat};
    \addplot[figthkq, forget plot] table[x=idx, y=kq_bool] {figures/data/threshold-min-df.dat};
    \addplot[figthssm, forget plot] table[x=idx, y=ssm_bool] {figures/data/threshold-min-df.dat};
  \end{axis}
  \end{tikzpicture}
  \caption{min-df：按文档频率下限裁词}
\end{subfigure}

\vspace{0.15em}
% ---------------- (c) min-match ----------------
\begin{subfigure}{0.49\linewidth}
  \centering
  \begin{tikzpicture}
  \begin{axis}[figthaxis,
    xlabel={min-match $k$}, ylabel={result documents},
    xtick={0,1,2}, xticklabels={1,2,3},
    xmin=-0.55, xmax=2.55,
  ]
    % 只有布尔查询有 min-match 数据（词/短语查询没有 k-of-n 语义）
    \addplot[figthkq, forget plot] table[x=idx, y=kq_bool] {figures/data/threshold-min-match.dat};
    \addplot[figthssm, forget plot] table[x=idx, y=ssm_bool] {figures/data/threshold-min-match.dat};
  \end{axis}
  \end{tikzpicture}
  \caption{min-match：要求命中 $k$ 个查询词（仅布尔查询）}
\end{subfigure}\hfill
% ---------------- (d) top-x ----------------
\begin{subfigure}{0.49\linewidth}
  \centering
  \begin{tikzpicture}
  \begin{axis}[figthaxis,
    xlabel={top-x cap},
    xtick={0,1,2}, xticklabels={5,10,20},
    xmin=-0.55, xmax=2.55,
  ]
    \addplot[figthlove, forget plot] table[x=idx, y=love] {figures/data/threshold-top-x.dat};
    \addplot[figthsword, forget plot] table[x=idx, y=sword] {figures/data/threshold-top-x.dat};
    \addplot[figthphrase, forget plot] table[x=idx, y=kr_phrase] {figures/data/threshold-top-x.dat};
    \addplot[figthkq, forget plot] table[x=idx, y=kq_bool] {figures/data/threshold-top-x.dat};
    \addplot[figthssm, forget plot] table[x=idx, y=ssm_bool] {figures/data/threshold-top-x.dat};
  \end{axis}
  \end{tikzpicture}
  \caption{top-x：只截断结果条数}
\end{subfigure}

\par\smallskip
\pgfplotslegendfromname{figthresholdlegend}
\caption{四个阈值改变的是查询语义而非结果数量：max-df/min-df 裁掉查询词
（单词查询结果可在临界档归零，双词语句裁词后结果反而变多），min-match 改
布尔条件，top-x 只截断条数；短语查询对前两个参数不敏感。}%
\label{fig:threshold}
\end{figure}
。
\pgfplotsset{
  % 面板公共设置
  figthaxis/.style={
    width=7.10cm, height=2.90cm,
    ymin=0,
    xtick style={draw=black!40},
    tick label style={font=\scriptsize},
    label style={font=\small},
    title style={font=\small},
    grid=major, grid style={black!8},
    unbounded coords=discard,
  },
  % 每条查询一条曲线：颜色/线型/标记三者组合区分
  figthlove/.style={color=mseblue, line width=0.9pt, mark=*, mark size=1.1pt},
  figthsword/.style={color=mseblue, dashed, line width=0.9pt,
                     mark=square*, mark size=1.0pt},
  figthphrase/.style={color=msegray, line width=0.9pt,
                      mark=triangle*, mark size=1.3pt},
  figthkq/.style={color=msegray, dashed, line width=0.9pt,
                  mark=diamond*, mark size=1.1pt},
  figthssm/.style={color=black!65, densely dotted, line width=1.0pt,
                   mark=o, mark size=1.1pt},
}
\begin{figure}[!htbp]
\centering
\captionsetup[subfigure]{font=scriptsize,skip=2pt}
% ---------------- (a) max-df ----------------
\begin{subfigure}{0.49\linewidth}
  \centering
  \begin{tikzpicture}
  \begin{axis}[figthaxis,
    xlabel={max-df threshold}, ylabel={result documents},
    xtick={0,1,2,3,4,5}, xticklabels={0.1,0.3,0.5,0.7,0.9,off},
    xmin=-0.55, xmax=5.55,
    legend to name={figthresholdlegend}, legend columns=3,
    legend style={draw=none, font=\scriptsize},
  ]
    \addplot[figthlove] table[x=idx, y=love] {figures/data/threshold-max-df.dat};
    \addlegendentry{love (word)}
    \addplot[figthsword] table[x=idx, y=sword] {figures/data/threshold-max-df.dat};
    \addlegendentry{sword (word)}
    \addplot[figthphrase] table[x=idx, y=kr_phrase] {figures/data/threshold-max-df.dat};
    \addlegendentry{king richard (phrase)}
    \addplot[figthkq] table[x=idx, y=kq_bool] {figures/data/threshold-max-df.dat};
    \addlegendentry{king queen (bool)}
    \addplot[figthssm] table[x=idx, y=ssm_bool] {figures/data/threshold-max-df.dat};
    \addlegendentry{sweet sword moon (bool)}
  \end{axis}
  \end{tikzpicture}
  \caption{max-df：裁掉 df 比例过高的查询词}
\end{subfigure}\hfill
% ---------------- (b) min-df ----------------
\begin{subfigure}{0.49\linewidth}
  \centering
  \begin{tikzpicture}
  \begin{axis}[figthaxis,
    xlabel={min-df threshold},
    xtick={0,1,2,3,4}, xticklabels={off,2,5,20,50},
    xmin=-0.55, xmax=4.55,
  ]
    \addplot[figthlove, forget plot] table[x=idx, y=love] {figures/data/threshold-min-df.dat};
    \addplot[figthsword, forget plot] table[x=idx, y=sword] {figures/data/threshold-min-df.dat};
    \addplot[figthphrase, forget plot] table[x=idx, y=kr_phrase] {figures/data/threshold-min-df.dat};
    \addplot[figthkq, forget plot] table[x=idx, y=kq_bool] {figures/data/threshold-min-df.dat};
    \addplot[figthssm, forget plot] table[x=idx, y=ssm_bool] {figures/data/threshold-min-df.dat};
  \end{axis}
  \end{tikzpicture}
  \caption{min-df：按文档频率下限裁词}
\end{subfigure}

\vspace{0.15em}
% ---------------- (c) min-match ----------------
\begin{subfigure}{0.49\linewidth}
  \centering
  \begin{tikzpicture}
  \begin{axis}[figthaxis,
    xlabel={min-match $k$}, ylabel={result documents},
    xtick={0,1,2}, xticklabels={1,2,3},
    xmin=-0.55, xmax=2.55,
  ]
    % 只有布尔查询有 min-match 数据（词/短语查询没有 k-of-n 语义）
    \addplot[figthkq, forget plot] table[x=idx, y=kq_bool] {figures/data/threshold-min-match.dat};
    \addplot[figthssm, forget plot] table[x=idx, y=ssm_bool] {figures/data/threshold-min-match.dat};
  \end{axis}
  \end{tikzpicture}
  \caption{min-match：要求命中 $k$ 个查询词（仅布尔查询）}
\end{subfigure}\hfill
% ---------------- (d) top-x ----------------
\begin{subfigure}{0.49\linewidth}
  \centering
  \begin{tikzpicture}
  \begin{axis}[figthaxis,
    xlabel={top-x cap},
    xtick={0,1,2}, xticklabels={5,10,20},
    xmin=-0.55, xmax=2.55,
  ]
    \addplot[figthlove, forget plot] table[x=idx, y=love] {figures/data/threshold-top-x.dat};
    \addplot[figthsword, forget plot] table[x=idx, y=sword] {figures/data/threshold-top-x.dat};
    \addplot[figthphrase, forget plot] table[x=idx, y=kr_phrase] {figures/data/threshold-top-x.dat};
    \addplot[figthkq, forget plot] table[x=idx, y=kq_bool] {figures/data/threshold-top-x.dat};
    \addplot[figthssm, forget plot] table[x=idx, y=ssm_bool] {figures/data/threshold-top-x.dat};
  \end{axis}
  \end{tikzpicture}
  \caption{top-x：只截断结果条数}
\end{subfigure}

\par\smallskip
\pgfplotslegendfromname{figthresholdlegend}
\caption{四个阈值改变的是查询语义而非结果数量：max-df/min-df 裁掉查询词
（单词查询结果可在临界档归零，双词语句裁词后结果反而变多），min-match 改
布尔条件，top-x 只截断条数；短语查询对前两个参数不敏感。}%
\label{fig:threshold}
\end{figure}
。
\pgfplotsset{
  % 面板公共设置
  figthaxis/.style={
    width=7.10cm, height=2.90cm,
    ymin=0,
    xtick style={draw=black!40},
    tick label style={font=\scriptsize},
    label style={font=\small},
    title style={font=\small},
    grid=major, grid style={black!8},
    unbounded coords=discard,
  },
  % 每条查询一条曲线：颜色/线型/标记三者组合区分
  figthlove/.style={color=mseblue, line width=0.9pt, mark=*, mark size=1.1pt},
  figthsword/.style={color=mseblue, dashed, line width=0.9pt,
                     mark=square*, mark size=1.0pt},
  figthphrase/.style={color=msegray, line width=0.9pt,
                      mark=triangle*, mark size=1.3pt},
  figthkq/.style={color=msegray, dashed, line width=0.9pt,
                  mark=diamond*, mark size=1.1pt},
  figthssm/.style={color=black!65, densely dotted, line width=1.0pt,
                   mark=o, mark size=1.1pt},
}
\begin{figure}[!htbp]
\centering
\captionsetup[subfigure]{font=scriptsize,skip=2pt}
% ---------------- (a) max-df ----------------
\begin{subfigure}{0.49\linewidth}
  \centering
  \begin{tikzpicture}
  \begin{axis}[figthaxis,
    xlabel={max-df threshold}, ylabel={result documents},
    xtick={0,1,2,3,4,5}, xticklabels={0.1,0.3,0.5,0.7,0.9,off},
    xmin=-0.55, xmax=5.55,
    legend to name={figthresholdlegend}, legend columns=3,
    legend style={draw=none, font=\scriptsize},
  ]
    \addplot[figthlove] table[x=idx, y=love] {figures/data/threshold-max-df.dat};
    \addlegendentry{love (word)}
    \addplot[figthsword] table[x=idx, y=sword] {figures/data/threshold-max-df.dat};
    \addlegendentry{sword (word)}
    \addplot[figthphrase] table[x=idx, y=kr_phrase] {figures/data/threshold-max-df.dat};
    \addlegendentry{king richard (phrase)}
    \addplot[figthkq] table[x=idx, y=kq_bool] {figures/data/threshold-max-df.dat};
    \addlegendentry{king queen (bool)}
    \addplot[figthssm] table[x=idx, y=ssm_bool] {figures/data/threshold-max-df.dat};
    \addlegendentry{sweet sword moon (bool)}
  \end{axis}
  \end{tikzpicture}
  \caption{max-df：裁掉 df 比例过高的查询词}
\end{subfigure}\hfill
% ---------------- (b) min-df ----------------
\begin{subfigure}{0.49\linewidth}
  \centering
  \begin{tikzpicture}
  \begin{axis}[figthaxis,
    xlabel={min-df threshold},
    xtick={0,1,2,3,4}, xticklabels={off,2,5,20,50},
    xmin=-0.55, xmax=4.55,
  ]
    \addplot[figthlove, forget plot] table[x=idx, y=love] {figures/data/threshold-min-df.dat};
    \addplot[figthsword, forget plot] table[x=idx, y=sword] {figures/data/threshold-min-df.dat};
    \addplot[figthphrase, forget plot] table[x=idx, y=kr_phrase] {figures/data/threshold-min-df.dat};
    \addplot[figthkq, forget plot] table[x=idx, y=kq_bool] {figures/data/threshold-min-df.dat};
    \addplot[figthssm, forget plot] table[x=idx, y=ssm_bool] {figures/data/threshold-min-df.dat};
  \end{axis}
  \end{tikzpicture}
  \caption{min-df：按文档频率下限裁词}
\end{subfigure}

\vspace{0.15em}
% ---------------- (c) min-match ----------------
\begin{subfigure}{0.49\linewidth}
  \centering
  \begin{tikzpicture}
  \begin{axis}[figthaxis,
    xlabel={min-match $k$}, ylabel={result documents},
    xtick={0,1,2}, xticklabels={1,2,3},
    xmin=-0.55, xmax=2.55,
  ]
    % 只有布尔查询有 min-match 数据（词/短语查询没有 k-of-n 语义）
    \addplot[figthkq, forget plot] table[x=idx, y=kq_bool] {figures/data/threshold-min-match.dat};
    \addplot[figthssm, forget plot] table[x=idx, y=ssm_bool] {figures/data/threshold-min-match.dat};
  \end{axis}
  \end{tikzpicture}
  \caption{min-match：要求命中 $k$ 个查询词（仅布尔查询）}
\end{subfigure}\hfill
% ---------------- (d) top-x ----------------
\begin{subfigure}{0.49\linewidth}
  \centering
  \begin{tikzpicture}
  \begin{axis}[figthaxis,
    xlabel={top-x cap},
    xtick={0,1,2}, xticklabels={5,10,20},
    xmin=-0.55, xmax=2.55,
  ]
    \addplot[figthlove, forget plot] table[x=idx, y=love] {figures/data/threshold-top-x.dat};
    \addplot[figthsword, forget plot] table[x=idx, y=sword] {figures/data/threshold-top-x.dat};
    \addplot[figthphrase, forget plot] table[x=idx, y=kr_phrase] {figures/data/threshold-top-x.dat};
    \addplot[figthkq, forget plot] table[x=idx, y=kq_bool] {figures/data/threshold-top-x.dat};
    \addplot[figthssm, forget plot] table[x=idx, y=ssm_bool] {figures/data/threshold-top-x.dat};
  \end{axis}
  \end{tikzpicture}
  \caption{top-x：只截断结果条数}
\end{subfigure}

\par\smallskip
\pgfplotslegendfromname{figthresholdlegend}
\caption{四个阈值改变的是查询语义而非结果数量：max-df/min-df 裁掉查询词
（单词查询结果可在临界档归零，双词语句裁词后结果反而变多），min-match 改
布尔条件，top-x 只截断条数；短语查询对前两个参数不敏感。}%
\label{fig:threshold}
\end{figure}


% src: tasks F37、F38、F39
表~\ref{tab:threshold-rows} 按参数归并代表性读数。
\begin{table}[!htbp]
  \centering
  \caption{四个查询阈值对返回结果的影响}
  \label{tab:threshold-rows}
  \footnotesize
  \renewcommand{\arraystretch}{1.0}
  \begin{tabular}{l l l c}
    \toprule
    参数 & 代表查询 & 取值：基线 $\to$ 结果 & Jaccard \\
    \midrule
    \code{--max-df} & \code{love} & 0.9/0.7：138 $\to$ 138/0 & 1.00/0.00 \\
    \code{--min-df} & \code{sweet sword moon} & 50：30 $\to$ 90 & 0.3333 \\
    \code{--min-match} &
      \makecell[l]{\code{king queen}\\\code{sweet sword moon}} &
      \makecell[l]{1/2/3：33 $\to$ 43/33/33\\1/2/3：30 $\to$ 93/39/30} &
      \makecell[c]{0.7674/1.00/1.00\\0.3226/0.7692/1.00} \\
    \code{--top-x} & \code{love} & 5/10/20：138 $\to$ 5/10/20 &
      0.0362/0.0725/0.1449 \\
    \bottomrule
  \end{tabular}
  \par\smallskip
  \footnotesize 注：Jaccard 为结果文档集合与基线集合的交并比；两个空集视为相同。  % src: tasks F38、F39
\end{table}

最高文档频率和最低文档频率筛选的是查询词，不是命中的文档。单词查询只有一个词项，
所以词项被删后结果会直接归零。love 的文档频率比值为 0.7077，在 0.7 档返回数从
138 降到 0；sweet 的比值为 0.4615，在 0.3 档也降到 0。

多词查询会先失去一项约束。sweet sword moon 的基线是 30 篇，使用
\code{--min-df 50} 后，文档频率为 37 和 35 的 sword、moon 被删，只剩 sweet，
于是三词全部命中退化为单词查询，结果反而增加到 90 篇。这个变化不能解释为“阈值提高了
精度”，它说明词项筛选已经改变了查询本身的含义。

最少命中词数在结果侧控制布尔条件。$k=1$ 相当于并集，$k=2$ 要求至少命中两词，
$k=3$ 对应三词全部命中；当 $k$ 超过可用词数时，程序将其限制为词项数。结果截断的
行为更直接：排序完成后只保留前 $x$ 篇，因此 love 的结果依次变为 5、10、20 篇，
不会改变返回文档之间的相对顺序。

短语和停用词给出了两类边界情况。短语查询绕过最大、最小文档频率阈值，king richard
在所有档上都保持 7 篇；如果删掉其中任何一个词，原短语就不再成立。停用词没有进入
索引，所以 the、of 在任何阈值下都返回 0。201 条测量中有 23 条 Jaccard 为 0，
逐条核对后，它们都发生在所有有效查询词项被删空之后；说明信息也记录了对应词项和
丢弃原因，没有出现无说明的空结果。

\paragraph{结论} 四个阈值并不是同一种“收紧”。最大、最小文档频率改变参与匹配的词项，
最少命中词数改变布尔条件，结果截断只改变输出范围。词项被删后，结果数可能归零，
也可能因为约束被削弱而上升。

\subsubsection{复杂度分析}
\label{sec:complexity}

表~\ref{tab:complexity} 给出主要操作的时间与空间上界。$N$、$T$、$P$、$D$、$H$ 分别
表示词元数、词项数、倒排记录数、文档数和命中文档数；$k$ 是查询词数，$R$ 是临时块数，
$B$ 是块预算，$C$ 是词典比较成本。短语验证中的 $\mathrm{tf}_{\max}$ 表示候选词项
在单个文档内的最大词频。

\begin{table}[!htbp]
  \centering
  \caption{主要操作的时间与空间复杂度上界}
  \label{tab:complexity}
  \footnotesize
  \begin{tabular}{l l l}
    \toprule
    操作 & 时间 & 空间 \\
    \midrule
    建索引（内存） & $O(N+\sum \mathrm{df}\cdot C)$ & $O(T+P+N)$ \\
    建索引（块式） & $O(N+RB\log B+TR)$ &
      $O(B+\text{当前块倒排记录})$ \\
    单词查询 & $O(\log T+\mathrm{df}+H\log H)$ & $O(D)$ \\
    短语查询 &
      $O(\mathrm{df}_{\min}\,\mathrm{tf}_{\max}\,k\log\mathrm{tf}_{\max}+H\log H)$ &
      $O(D+k)$ \\
    布尔查询 & $O(\sum \mathrm{df}+H\log H)$ & $O(D)$ \\
    第二版索引编码 & $O(P+N)$ & $O(\text{单块})$ \\
    B+ 树查找 & $O(h\log F)$，另需 $O(h)$ 次页读 & $O(h)$ 页缓存 \\
    \bottomrule
  \end{tabular}
\end{table}

内存建索引的时间主要花在每个词元的词典访问上，输出部分则同时保存词项表、倒排记录和
位置。三种词典写出的索引大小相同，改变的只是比较次数和比较常数。块式构建用 $B$
限制同时驻留的活词，但不能消掉全部成本：$B$ 越小，临时块越多，排序和归并的代价
越容易显现。磁盘 B+ 树换了一种取舍，查询只读取根到叶的 $O(h)$ 个页，不再要求完整
词项表常驻内存。

查询侧的主要成本是遍历相关倒排记录，以及排序命中文档。短语查询先选文档频率最小的
词项，可以缩小候选集，但仍要对每个候选验证 $k$ 个位置。阈值与复杂度的联系也很直接：
最大、最小文档频率在遍历倒排记录之前删词，可能减少后续工作；最少命中词数和结果截断
发生在计分之后，其中 \code{--top-x} 只减少输出数量，不会省掉扫描和排序。对于当前
实现，阈值首先是改变查询语义的工具，不是统一的性能优化开关。

\setlength{\textfloatsep}{20pt plus 2pt minus 4pt}
\setlength{\floatsep}{12pt plus 2pt minus 2pt}
\setlength{\intextsep}{12pt plus 2pt minus 2pt}
